\documentclass[11pt]{article}
\usepackage[a4paper,margin=2.2cm]{geometry}
\usepackage{amsmath,amssymb}
\usepackage{booktabs}
\usepackage{xcolor}
\usepackage{hyperref}
\usepackage{enumitem}

\newcommand{\meaningpar}[1]{\par\noindent\textit{Meaning of the regressors.} #1}

\title{\vspace{-1.5em}MetaFVG Trend-Filter Family: Regression \& Lasso Formulations}
\author{Research notes --- \texttt{metalib/research/}}
\date{Reference setup: HTF=4h, LTF=M15, post lookahead-bias fix}

\begin{document}
\maketitle
\vspace{-2em}

\section*{Notation}
Let $c_1,\dots,c_n$ denote the HTF (4h) close series and, at bar $i$, let $y = (c_{i-w+1},\dots,c_i)$
be the trailing window of width $w$ (default $w=20$), regressed against the bar-index
regressor $x = (0,1,\dots,w-1)$. All quantities below are computed causally: the value
attributed to bar $i$ uses only $c_{\le i}$.
\meaningpar{$x$ is \emph{not} calendar time --- it is a purely local index that resets to
$0$ at the start of every window and counts up to $w-1$. Every rolling fit is therefore
``how has price moved over the last $w$ bars,'' never an absolute-time trend. $d\in\{+1,-1\}$
throughout denotes the candidate trade direction already implied by the HTF FVG zone the
current price sits in (from the signal pipeline, not from the regression itself) --- every
gate below only asks whether the regression \emph{agrees} with that direction.}

\section{Regression Gate}
A rolling univariate OLS of price on bar index within each trailing window:
\[
c_t = \beta_0 + \beta_1 x_t + \varepsilon_t, \qquad
\hat\beta_1 = \frac{\sum_t (x_t-\bar x)(y_t-\bar y)}{\sum_t (x_t-\bar x)^2}, \qquad
\hat\beta_0 = \bar y - \hat\beta_1 \bar x
\]
\[
R^2 = 1 - \frac{\mathrm{SS}_{\mathrm{res}}}{\mathrm{SS}_{\mathrm{tot}}}, \qquad
\mathrm{SS}_{\mathrm{res}} = \sum_t (y_t - \hat y_t)^2, \qquad
\mathrm{SS}_{\mathrm{tot}} = \sum_t (y_t - \bar y)^2
\]
\textbf{Gate condition} (a signal in direction $d\in\{+1,-1\}$ is accepted iff):
\[
\operatorname{sign}(\hat\beta_1) = d \quad \text{and} \quad R^2 \ge \tau_{R^2}, \qquad \tau_{R^2}=0.5 \text{ (default)}
\]
\meaningpar{$c_t=y_t$ is simply the price at bar $t$ within the window --- the thing being explained.
$x_t$ is that bar's position in the window (see Notation): the regression asks ``as $x$ counts
up from $0$ to $w-1$, does price move consistently in one direction?'' $\hat\beta_1$ (the slope) is
the fitted \emph{price change per bar} --- its sign is the trend direction, its magnitude is
irrelevant to the gate (only used to determine sign). $R^2\in[0,1]$ is the fraction of the window's
price variance explained by that straight line --- a ``cleanliness'' score: $R^2$ near $1$ means
price moved almost monotonically, near $0$ means it was choppy/range-bound even if the fitted
slope happens to point the right way. $\tau_{R^2}=0.5$ is the minimum cleanliness required before
the gate trusts the direction call at all. This is the config that survived the HTF lookahead-bias
fix with a real, if modest, edge (avg.\ individual Sharpe $+0.010$, equal-weighted 7-pair FX
Majors portfolio Sharpe $+0.143$).}

\section{Theil-Sen Gate}
Identical $R^2$ gate to \S1, but the direction call uses the Theil--Sen (median-of-pairwise-slopes)
estimator instead of the OLS slope --- robust to a single outlier bar dominating the fit:
\[
\hat\beta_1^{\mathrm{TS}} = \operatorname{median}\left\{ \frac{y_j - y_i}{j - i} \;:\; 0\le i<j< w \right\}
\]
\[
\text{Gate: } \operatorname{sign}(\hat\beta_1^{\mathrm{TS}}) = d \quad \text{and} \quad R^2 \ge \tau_{R^2}
\]
\meaningpar{For every pair of bars $i<j$ in the window, $(y_j-y_i)/(j-i)$ is the average
price-change-per-bar \emph{measured only between those two bars} --- i.e.\ ``if price only
moved between these two points, what would the slope be.'' Taking the median over all
$\binom{w}{2}$ such pairs means one freak bar (a single large spike) can shift the result by
at most one rank, unlike OLS where it can dominate the whole fit through the squared-error
term. $R^2$ here is still the \emph{OLS} goodness-of-fit from \S1 --- only the direction call
changes, the cleanliness gate does not. Empirically tied Regression Gate almost exactly (avg.\
Sharpe $+0.072$ vs.\ $+0.073$ pre-refinement scale), i.e.\ outlier bars were not materially
flipping the OLS direction call on this asset set.}

\section{Spearman Gate}
Replaces both the OLS slope \emph{and} $R^2$ with a single rank-correlation statistic, testing
monotonicity rather than linearity (a nonlinear but consistently directional move is not penalized
for curvature the way $R^2$ penalizes it):
\[
\rho = 1 - \frac{6\sum_t d_t^2}{w(w^2-1)}, \qquad d_t = \operatorname{rank}(y_t) - t
\]
\[
\text{Gate: } \operatorname{sign}(\rho) = d \quad \text{and} \quad |\rho| \ge \tau_\rho
\]
\meaningpar{$\operatorname{rank}(y_t)$ is $y_t$'s rank (1st-smallest, 2nd-smallest, \dots) among
the $w$ prices in the window; $t$ itself is already the bar's rank in time (since $x=0,\dots,w-1$
is already sorted). $d_t$ is therefore the gap between ``where does this bar rank by price'' and
``where does this bar rank by time'' --- if price rose monotonically, the two rankings are
identical and every $d_t=0$. $\rho\in[-1,1]$ (Spearman's rank correlation) summarizes that
agreement in one number: $+1$ means price and time are perfectly rank-ordered together
(monotonic uptrend, any shape), $-1$ means perfectly monotonic down, $0$ means no consistent
ordering. Unlike $R^2$, $\rho$ does not care whether the monotonic move was a straight line or a
curve --- only that it never reversed rank. Tested at $\tau_\rho \in \{0.5,\ 0.82,\ 0.85\}$; none
beat Regression Gate on FX Majors post-fix (avg.\ Sharpe $-0.011$ to $-0.254$ depending on
threshold and whether risk-sizing is layered on top).}

\section{Lasso Trend Gate (walk-forward)}
A genuinely predictive, out-of-sample model: fit on a trailing window of \texttt{train\_bars}
(450) bars, refit every \texttt{test\_bars} (190) bars, predicting the forward $k$-bar
($k=5$) log return:
\[
Y_t = \ln c_{t+k} - \ln c_t
\]
Regressor vector at bar $t$ (base feature set, before EMA augmentation):
\[
X_t = \big[\,\mathrm{slope}_t,\ R^2_t,\ \sigma^{\mathrm{short}}_t,\ \sigma^{\mathrm{long}}_t,\
\rho^{\mathrm{vol}}_t,\ \mathrm{slope}_t\rho^{\mathrm{vol}}_t,\ R^2_t\rho^{\mathrm{vol}}_t,\
\mathrm{slope}_t\sigma^{\mathrm{short}}_t,\ 1 \,\big]
\]
where $\sigma^{\mathrm{short}}_t,\sigma^{\mathrm{long}}_t$ are realized volatility of log returns
over 10- and 60-bar trailing windows, and $\rho^{\mathrm{vol}}_t = \sigma^{\mathrm{short}}_t/\sigma^{\mathrm{long}}_t$.
Fit by $L_1$-penalized least squares (numba coordinate descent --- this environment's BLAS crashes
on any matrix-matrix multiply, so no \texttt{sklearn}/\texttt{statsmodels} closed-form solve):
\[
\hat\beta = \operatorname*{arg\,min}_\beta \ \frac{1}{2n}\lVert Y - X\beta \rVert_2^2 + \alpha \lVert \beta \rVert_1,
\qquad \alpha \text{ chosen by 5-fold CV over 60 candidate values}
\]
\textbf{Gate condition}, with $\hat Y_t = X_t^\top\hat\beta$ (clipped to the training target range)
and $\tau_t$ a \emph{causal} rolling quantile of $|\hat Y|$ shifted one bar (never uses a
prediction generated at or after $t$):
\[
\operatorname{sign}(\hat Y_t) = d \quad \text{and} \quad |\hat Y_t| \ge \tau_t
\]
\meaningpar{$Y_t$, the \emph{target} being predicted, is ``the log return you would realize if you
held from now ($t$) to $k=5$ HTF bars later'' --- this is the one genuinely forward-looking
quantity in the whole document, and it is only ever used as a training label, computed from
\emph{fully elapsed} history at fit time (never a live, in-progress bar).
Reading $X_t$ component by component: $\mathrm{slope}_t,R^2_t$ are exactly \S1's OLS
outputs, but here they are \emph{inputs} to a model rather than a hard pass/fail gate --- the
Lasso decides how much weight (if any) they deserve. $\sigma^{\mathrm{short}}_t$ (10-bar) and
$\sigma^{\mathrm{long}}_t$ (60-bar) are realized volatility of log returns --- ``how much has
price been jumping around lately,'' measured over a quick and a slower window respectively.
$\rho^{\mathrm{vol}}_t=\sigma^{\mathrm{short}}_t/\sigma^{\mathrm{long}}_t$ is a vol-\emph{regime}
signal: $>1$ means volatility is currently expanding relative to its own recent baseline, $<1$
means it is contracting. The three products
($\mathrm{slope}_t\rho^{\mathrm{vol}}_t$, $R^2_t\rho^{\mathrm{vol}}_t$,
$\mathrm{slope}_t\sigma^{\mathrm{short}}_t$) are curated interaction terms, e.g.\ ``does a strong
trend behave differently when volatility is expanding vs.\ contracting'' --- a small, hand-picked
set rather than a full polynomial expansion, so the model stays interpretable by inspection of
which coefficients the Lasso keeps nonzero. The trailing constant $1$ is the intercept regressor.
$\hat\beta$ is the fitted coefficient vector: the $L_1$ penalty shrinks the coefficient on any
feature that doesn't earn its keep exactly to zero, so $\hat\beta$ itself is a readable statement
of ``which of these nine things actually mattered this refit cycle.'' $\alpha$ controls how harshly
that shrinkage is applied (larger $\alpha$ $\to$ fewer, smaller coefficients, chosen by
cross-validation rather than fixed by hand). $\hat Y_t=X_t^\top\hat\beta$ is the model's forecast
of $Y_t$ using only information available at $t$; $\tau_t$ is a magnitude bar on that forecast ---
only trust predictions that are unusually large relative to the model's own recent prediction
history, not just any nonzero sign.}

\section{EMA-Augmented Lasso --- Price-Based (superseded)}
First attempt at adding EMAs as regressors: normalized price-vs-EMA distance and crossover,
appended to $X_t$ above, with $\mathrm{EMA}_9,\mathrm{EMA}_{21}$ computed on \emph{price}:
\[
\mathrm{ema\_fast\_dist}_t = \frac{c_t - \mathrm{EMA}_9(c)_t}{c_t}, \qquad
\mathrm{ema\_slow\_dist}_t = \frac{c_t - \mathrm{EMA}_{21}(c)_t}{c_t}
\]
\[
\mathrm{ema\_cross}_t = \frac{\mathrm{EMA}_9(c)_t - \mathrm{EMA}_{21}(c)_t}{c_t}, \qquad
\text{plus interaction } \mathrm{ema\_cross}_t \cdot \rho^{\mathrm{vol}}_t
\]
\meaningpar{$\mathrm{EMA}_9(c)$ and $\mathrm{EMA}_{21}(c)$ are exponentially-smoothed price ---
a ``fast'' running average that reacts within about 9 bars and a ``slow'' one that reacts over
about 21 (the classic short/long CTA pairing; see \S6 for the exact recursion). Dividing by
$c_t$ turns each raw price-minus-average into a \emph{normalized} distance so it is comparable
across instruments/price scales: $\mathrm{ema\_fast\_dist}_t>0$ means price sits above its own
short-run average (recently strong), and similarly for the slow one against the longer-run
average. $\mathrm{ema\_cross}_t$ is the gap between the fast and slow average itself, again
normalized by price --- positive when the fast average has pulled above the slow one, the
textbook ``golden cross'' direction signal. Underperformed Regression Gate at both tested
thresholds (avg.\ Sharpe $-0.308$ loose, $-0.257$ tight; portfolio blend \emph{worse} than the
individual average in both cases, unlike Regression Gate where blending helps).}

\section{EMA-Augmented Lasso --- Vol-Adjusted Returns (current)}
Per your request: the EMA is taken of \emph{vol-adjusted returns}, not price, so the smoothed
series is comparable across volatility regimes:
\[
r_t = \ln c_t - \ln c_{t-1}, \qquad
r^{\mathrm{vadj}}_t = \frac{r_t}{\sigma^{\mathrm{short}}_t}
\]
\[
\mathrm{ema\_fast}_t = \mathrm{EMA}_9(r^{\mathrm{vadj}})_t, \qquad
\mathrm{ema\_slow}_t = \mathrm{EMA}_{21}(r^{\mathrm{vadj}})_t
\]
\[
\mathrm{ema\_cross}^{\mathrm{vadj}}_t = \mathrm{ema\_fast}_t - \mathrm{ema\_slow}_t, \qquad
\text{plus interaction } \mathrm{ema\_cross}^{\mathrm{vadj}}_t \cdot \rho^{\mathrm{vol}}_t
\]
where, for a series $z$, $\mathrm{EMA}_p(z)_t = w z_t + (1-w)\mathrm{EMA}_p(z)_{t-1}$,
$w = 2/(p+1)$, seeded from the first non-NaN value and held steady across any NaN gap.
\meaningpar{$r_t$ is the plain bar-to-bar log return --- ``how much did price move this bar,''
with no normalization. $r^{\mathrm{vadj}}_t$ divides that by $\sigma^{\mathrm{short}}_t$ (the same
10-bar realized vol from \S4), turning it into ``how many typical-recent-moves was this bar's
return,'' roughly a rolling z-score of the return: a $0.5\%$ move counts as large during a calm
stretch and as unremarkable during a volatile one, whereas a raw-price EMA cannot distinguish the
two. $\mathrm{ema\_fast}_t,\mathrm{ema\_slow}_t$ are no longer distances from anything --- since
the base series is already a (vol-adjusted) return, not a price level, the smoothed value itself
\emph{is} the signal: a running estimate of recent vol-adjusted momentum, fast (9-bar) and slow
(21-bar). $\mathrm{ema\_cross}^{\mathrm{vadj}}_t$ is their difference, i.e.\ momentum
\emph{acceleration} --- positive when recent momentum is running hotter than its own longer-run
average. In the smoothing recursion, $w=2/(p+1)$ is the EMA weight (not to be confused with the
regression window width $w$ used in \S1--\S3 --- unfortunately the same letter, distinguished by
context in the source). Recovers some ground versus the price-based version but still loses to
Regression Gate (avg.\ Sharpe $-0.239$ loose, $-0.230$ tight; blend still worse than the individual
average).}

\section{Shelved: LTF Significance-Gated OLS Bias}
From the (rejected) ``back to basics'' redesign: same closed-form OLS as \S1, but applied to the
\emph{LTF} (M15) close series, and gated on statistical significance of the slope rather than
goodness-of-fit:
\[
\mathrm{SE}(\hat\beta_1) = \sqrt{\frac{\hat\sigma^2}{\sum_t(x_t-\bar x)^2}}, \qquad
\hat\sigma^2 = \frac{\mathrm{SS}_{\mathrm{res}}}{w-2}, \qquad
t = \frac{\hat\beta_1}{\mathrm{SE}(\hat\beta_1)}
\]
\[
\text{Gate: } \operatorname{sign}(\hat\beta_1) = d \quad \text{and} \quad |t| \ge \tau_t
\]
\meaningpar{$\hat\sigma^2$ is the fit's residual variance --- how noisy the window is around its
own best-fit line, degrees-of-freedom-corrected. $\mathrm{SE}(\hat\beta_1)$ turns that into an
uncertainty estimate specifically for the slope: how much would $\hat\beta_1$ plausibly wobble if
re-fit on similar noise. The $t$-statistic is the slope measured in units of its own uncertainty
--- ``is this slope big compared to how unsure we are about it,'' as opposed to $R^2$ in \S1 which
only asks ``does a line fit well'' without any notion of estimation uncertainty. $\tau_t$ is the
minimum $|t|$ required to trust the direction call, analogous to $\tau_{R^2}$ in \S1 but on a
statistical-significance scale instead of a goodness-of-fit scale. Tested $\tau_t \in [0.5, 4.0]$
across all 7 FX Majors: average Sharpe was negative throughout and \emph{worsened} as $\tau_t$
tightened ($-0.286$ at $\tau_t{=}1.5$ down to $-0.370$ at $\tau_t{=}3.5$) --- the opposite of the
single-symbol (EURUSD) read that motivated testing it, and GBPUSD flipped from this campaign's
best individual result (Sharpe $1.10$ under \S1) to one of this variant's worst ($-1.03$). Shelved;
code retained in \texttt{metafvg\_basics\_variant.py} for reference.}

\section*{Summary}
\begin{center}
\begin{tabular}{lcc}
\toprule
Config & Avg.\ individual Sharpe & Blended portfolio Sharpe \\
\midrule
Regression Gate & $+0.010$ & $+0.143$ \\
Theil-Sen Gate & $+0.072$\textsuperscript{*} & --- \\
Spearman Gate (q=0.5) & $-0.011$\textsuperscript{*} & --- \\
Lasso Trend Gate (base features) & $-0.16$ to $-0.42$\textsuperscript{*} & --- \\
EMA-Augmented Lasso (price) & $-0.308$ / $-0.257$ & $-0.688$ / $-0.712$ \\
EMA-Augmented Lasso (vol-adj.) & $-0.239$ / $-0.230$ & $-0.488$ / $-0.640$ \\
\bottomrule
\end{tabular}
\end{center}
{\small\textsuperscript{*}Full-universe (47-symbol) scale, not the FX-Majors-only re-verification
used for the rows below it; included for context, not a like-for-like comparison.}

\end{document}
